\scanpage{81}
% loc 1, top heading and date
\uline{Azumaya rang} 4 \mspage{I} Avril \unsure{85|81|86}.

% loc 2–4
1. Soient $k$ anneau, $A$ algèbre d'Azumaya sur $k$ de rang 4 (donc loc. iso. à $M_2(k)$). On a alors
% loc 5–6
\[
\left\{\begin{aligned}
\operatorname{Tr}: A&\longrightarrow k & k\text{-linéaire}\\
\det: A&\longrightarrow k & \text{multiplicatif, quadratique}
\end{aligned}\right. \tag*{(1)}
\]
% loc 7
tels que $\forall u\in A$, on ait
% loc 8
\[ u^2-\operatorname{Tr}(u).u+\det(u)=1 \tag*{(2)} \]
% loc 9
(Hamilton-Cayley). De plus, on a
% loc 10
\[ \det(u+v)=\det u+\det v+(\operatorname{Tr}u\operatorname{Tr}v-\operatorname{Tr}(uv)) \tag*{(3)} \]
% loc 11
(i.e.
% loc 12
\[ (u,v)\longmapsto\unsure{\psi|\Psi|\Phi}(u,v)\overset{\text{dfn}}{=}\operatorname{Tr}u\operatorname{Tr}v-\operatorname{Tr}(uv) \]
% loc 13
est la forme bil. \unsure{ass.} à la forme quad. $\det$).

% loc 14–19, remark
\uline{\unsure{Remarque|Remarques}} On peut s'intéresser aux $k$-Alg. quelconques $A$ pour lesquelles on a (1) (2) (3). Si $A$ est proj. de t.f. disons et 2 inv. dans $k$, alors $(1)(2)\Longrightarrow(3)$, car il suffit de le voir pour $u=v$ (deux fonctions bil. \uline{sym} qui coïncident pour $u=v$ coïncident)\unsure{.|,}
% loc 19–22, cancelled continuation and restart
\deleted{et alors cela se \inserted{résulte de} réduit à une vérification dans les $k[u]=k\oplus k.u$ quadratiques sur $k$ \unsure{...})} La relation pertinente est
% loc 23
\[ (\operatorname{Tr}u)^2-\operatorname{Tr}(u^2)=2\det u, \tag*{(4)} \]
% loc 24
et elle résulte de (2) en \unsure{s'écrisant|l'écrivant} sous forme
% loc 25, between bottom lines
\[ u^2=\operatorname{Tr}(u)u-\det(u) \tag*{(2')} \]
% loc 26, bottom edge
et portant cette valeur de $u^2$ dans (4), compte tenu
% archive sheet number: 80
